\section{Related Work}

\subsection{Multimodal Semantic Segmentation}

RGB-X semantic segmentation improves robustness by incorporating thermal, depth, polarization, LiDAR, or other auxiliary modalities, and comprehensive overviews can be found in prior surveys~\cite{zhang2021multimodalsurvey}. Early RGB-T methods mainly relied on dual-branch CNN encoders and explicit fusion decoders, as exemplified by MFNet, RTFNet, and FuseNet~\cite{ha2017mfnet,sun2019rtfnet,hazirbas2016fusenet}. Later work explored stronger feature interaction through self-supervised adaptation, non-local fusion, attention-based fusion, complementarity-aware modeling, and real-time multimodal decoding~\cite{valada2020selfsupervised,yan2021nlfnet,xu2021attentionfusion,wu2022complementarity,sun2020realtimefusion}. More recent methods, including GMNet, CMX, CMNeXt, Multimodal Token Fusion, GeminiFusion, RoadFormer, RoadFormer+, and CAFuser, move further toward transformer-based or condition-aware cross-modal interaction~\cite{zhou2021gmnet,zhang2023cmx,zhang2023deliver,wang2022multimodaltoken,jia2024geminifusion,li2024roadformer,huang2025roadformerplus,broedermann2025cafuser}. Despite this progress, most of these methods focus on driving or generic road scenes rather than drone-view perception with large scale variation and severe modality-quality fluctuations.

\subsection{Foundation Model Adaptation for Dense Prediction}

The success of large-scale visual pre-training has made foundation-model adaptation a central topic in dense prediction. ViT and MAE established strong transformer-based visual backbones, while ConvNeXt showed that modernized ConvNets can remain highly competitive~\cite{dosovitskiy2021vit,he2022mae,liu2022convnext}. ViT-Adapter demonstrated that dense prediction performance can be substantially improved by attaching lightweight trainable modules to frozen pre-trained transformers~\cite{chen2023vitadapter}. In parallel, parameter-efficient tuning methods such as LoRA, Conv-LoRA, and recent PEFT surveys highlight the practical value of updating only a small set of trainable parameters when adapting large pre-trained models~\cite{hu2022lora,zhong2024convlora,xin2024peftsurvey}. For segmentation foundation models, SAM-Adapter and SAM2-Adapter extended this philosophy to SAM and SAM~2, while MM SAM-Adapter further injected auxiliary-modal information into a SAM-based RGB backbone for multimodal semantic segmentation~\cite{chen2023samadapter,chen2024sam2adapter,curti2025mmsamadapter}. Our work follows this adaptation route, but targets drone-view RGB-T segmentation and emphasizes condition-aware fusion and small-object-aware decoding.

\subsection{Drone-View Scene Parsing}

Compared with ground-view RGB-T segmentation, drone-view segmentation involves stronger perspective distortion, broader scale variation, and much more severe small-object degradation. UAVid, Agriculture-Vision, and SynDrone illustrate these characteristics from different aerial or UAV scenarios~\cite{lyu2020uavid,chiu2020agriculturevision,rizzoli2023syndrone}. Architectures such as UAVformer and 2DSegFormer explicitly target aerial semantic segmentation with stronger global modeling~\cite{zhang2023uavformer,li2022twodsegformer}. On the remote-sensing side, RS-SAM, Multiscale Adapter, and SAM2Former investigate how SAM-family models can be adapted to high-resolution aerial and remote-sensing segmentation~\cite{zhang2024rssam,chen2025multiscaleadapter,li2025sam2former}. However, these methods mainly focus on RGB imagery or remote-sensing semantics rather than RGB-T fusion under drone-specific modality instability. Our design therefore focuses on two aspects that remain under-addressed in this setting: condition-dependent multimodal fusion and the recovery of small targets in hierarchical decoding.
